
\section{Hypocentral homomorphisms and factorization system}
\label{section: hypocentral homomorphisms}

In this section, we extend relative lower central series to transfinite ordinals and introduce hypocentral homomorphisms. The main result is an orthogonal factorization system (for surjections) whose left class consists of strongly connected homomorphisms.

\subsection{Hypocentral homomorphisms of quandles}
\label{subsection: hypocentral homomorphisms of quandles}
We introduce the notion of a \emph{hypocentral homomorphism} and first establish some basic properties parallel to those of nilpotent homomorphisms.

\begin{dfn}
\label[dfn]{definition: transfinite relative lower central series}
Let $N\trianglelefteq G$ be a normal subgroup. We define the \emph{transfinite $G$-relative lower central series} of $N$ in the following way: $\relLCS_0(G,N)\coloneqq N$ and $\relLCS_{\alpha+1}(G,N)\coloneqq[G,\relLCS_\alpha(G,N)]$ for an ordinal $\alpha$. For every non-zero limit ordinal $\lambda$, put
\[
    \relLCS_\lambda(G,N)
    \coloneqq
    \bigcap_{\alpha<\lambda}\relLCS_\alpha(G,N).
\]
We say $N$ is \emph{$G$-hypocentral} if $\relLCS_\alpha(G,N)=1$ for some ordinal $\alpha$.
\end{dfn}

\begin{lem}
\label[lem]{lemma: relative hypocenter}
    Let $N\trianglelefteq G$ be a normal subgroup. The transfinite $G$-relative lower central series of $N$ stabilizes, and its stable term is the largest $G$-perfect normal subgroup of $G$ contained in $N$.
\end{lem}

\begin{proof}
The series stabilizes since $G$ has only set many normal subgroups. Its stable term is $G$-perfect. If $K\trianglelefteq G$ satisfies $K\subseteq N$ and $[G,K]=K$, then transfinite induction gives $K\subseteq\relLCS_\alpha(G,N)$ for every ordinal $\alpha$. Thus, $K$ is contained in the stable term.
\end{proof}

We write $\relHypCenter{G}{N}$ for the stable term and call it the \emph{$G$-relative hypocenter} of $N$. Thus, $N$ is $G$-hypocentral if and only if $\relHypCenter{G}{N}=1$.

\begin{cor}
\label{corollary:relative_hypocentral_iff_no_non-trivial_relative_perfect}
    A normal subgroup $N\trianglelefteq G$ is $G$-hypocentral if and only if there is no non-trivial $G$-perfect normal subgroup of $G$ contained in $N$.
\end{cor}

In turn, we extend the notion of nilpotency to the transfinite setting.

\begin{dfn}
\label[dfn]{definition: hypocentral homomorphism}
    Let $f\colon P\twoheadrightarrow Q$ be a surjective quandle homomorphism. We call $f$ \emph{hypocentral} if $\relInn(f)$ is $\Inn(P)$-hypocentral. 
\end{dfn}

Every nilpotent homomorphism is hypocentral. For the unique homomorphism $P\twoheadrightarrow\{*\}$, the preceding definition is the usual hypocentrality condition for the group $\Inn(P)$.

For every ordinal $\alpha$, we use the notation $\relInnLCS{\alpha}{f}$ and $\relTransLCS{\alpha}{f}$ for the $\alpha$-th term of the transfinite relative lower central series of $\relInn(f)$ and $\relTrans(f)$ in $\Inn(P)$, respectively.

\begin{lem}
\label[lem]{lemma: transfinite interlacing of the two relative lower central series}
Let $f\colon P\twoheadrightarrow Q$ be a surjective quandle homomorphism. Then, for every ordinal $\alpha$, we have
\[ \relInnLCS{\alpha+1}{f} \subseteq \relTransLCS{\alpha}{f} \subseteq \relInnLCS{\alpha}{f}. \]
Moreover, $\relInnLCS{\alpha}{f}=\relTransLCS{\alpha}{f}$ for $\alpha\geq\omega$.
\end{lem}

\begin{proof}
Put $G\coloneqq\Inn(P)$, $N\coloneqq\relInn(f)$, and $T\coloneqq\relTrans(f)$. The second inclusion follows by transfinite induction from $T\subseteq N$. Since $[G,N]\subseteq T$ by \cref{lem:commutator_inn_rel_trans_rel}, the first inclusion holds for $\alpha=0$, and the successor step follows by taking commutators with $G$. If $\lambda$ is a limit ordinal, then $\relLCS_\lambda(G,N)\subseteq\relLCS_{\beta+1}(G,N) \subseteq\relLCS_\beta(G,T)$ for every $\beta<\lambda$. Hence, $[G,\relLCS_\lambda(G,N)]\subseteq\relLCS_\lambda(G,T)$.

Taking intersections over the finite stages gives $\relInnLCS{\omega}{f}=\relTransLCS{\omega}{f}$. The equality for $\alpha\geq\omega$ follows by transfinite induction.
\end{proof}

For a surjective quandle homomorphism $f$, \cref{lemma: transfinite interlacing of the two relative lower central series} shows that the stable terms of the two transfinite relative lower central series associated with $\relInn(f)$ and $\relTrans(f)$ coincide. We will denote this common term by $\textstyle \hypCenter{f}\coloneqq \bigcap_\alpha \relInnLCS{\alpha}{f} = \bigcap_{\alpha} \relTransLCS{\alpha}{f}$.

\begin{cor}
\label[cor]{corollary: relative hypocenter and hypocentral lengths}
A surjective homomorphism $f\colon P\twoheadrightarrow Q$ is hypocentral if and only if the normal subgroup $\relTrans(f)$ is $\Inn(P)$-hypocentral.
\end{cor}

\begin{proof}
This follows from \cref{lemma: transfinite interlacing of the two relative lower central series}.
\end{proof}

\begin{rem}
We call $f$ \emph{residually nilpotent} if $\textstyle \relInnLCS{\omega}{f}=\bigcap_{i<\omega}\relInnLCS{i}{f}=1$. Every residually nilpotent homomorphism is hypocentral.
\end{rem}




\subsection{Characterizations of hypocentral homomorphisms}
\label{subsection: characterizations of hypocentral homomorphisms}

In this subsection, we give characterizations of hypocentral homomorphisms in terms of $\relTrans(f)$, $\relAs(f)$, and $\Inn(P)$-perfect subgroups.
To establish these characterizations, we first record a group-theoretic observation used for the relative adjoint group.

\begin{lem}
\label[lem]{lemma: hypocentrality under central extensions}
    Let $\pi\colon A\twoheadrightarrow G$ be a central extension of groups, let $R\trianglelefteq A$ be a normal subgroup, and put $T\coloneqq\pi(R)$. Then, $R$ is $A$-hypocentral if and only if $T$ is $G$-hypocentral.
\end{lem}

\begin{proof}
($\Leftarrow$): Let $L\trianglelefteq A$ be an $A$-perfect subgroup contained in $R$. Then, since $[G,\pi(L)]=\pi([A,L])=\pi(L)$, the image $\pi(L)$ is $G$-perfect and contained in $T=\pi(R)$. By the assumption that $T$ is $G$-hypocentral, \cref{corollary:relative_hypocentral_iff_no_non-trivial_relative_perfect} yields $\pi(L)=1$. Therefore, we have $L\subseteq \Ker(\pi) \subseteq Z(A)$, which implies $L=[A,L]=1$. Thus $R$ is $A$-hypocentral by \cref{corollary:relative_hypocentral_iff_no_non-trivial_relative_perfect}.

($\Rightarrow$): Let $K\trianglelefteq G$ be a $G$-perfect subgroup contained in $T$. Put $\widetilde{K}\coloneqq(\pi\rvert_R)^{-1}(K) =R\cap\pi^{-1}(K)$ and $L\coloneqq[A,\widetilde{K}]$. Then, $\widetilde{K}\trianglelefteq A$, and hence $L\trianglelefteq A$ and $L\subseteq\widetilde{K}\subseteq R$. Moreover, $\pi(\widetilde{K})=K$ and $\pi(L)=[G,K]=K$.

We claim that $\widetilde{K} = L\bigl(\widetilde{K}\cap\Ker(\pi)\bigr)$. Indeed, the inclusion from right to left is clear. Conversely, let $x\in\widetilde{K}$. Since $\pi(L)=K=\pi(\widetilde{K})$, there is $l\in L$ such that $\pi(l)=\pi(x)$. Then $l^{-1}x\in\widetilde{K}\cap\Ker(\pi)$, and hence $x\in L(\widetilde{K}\cap\Ker(\pi))$.

Since $\widetilde{K}\cap\Ker(\pi)\subseteq Z(A)$, this equality gives
\[ L = [A,\widetilde K] = [A,L(\widetilde K\cap\Ker(\pi))] = [A,L]. \]
Thus, $L$ is an $A$-perfect subgroup contained in $R$, and hence must be trivial by \cref{corollary:relative_hypocentral_iff_no_non-trivial_relative_perfect}. It follows that $K=\pi(L)=1$, which shows that $T$ is $G$-hypocentral.
\end{proof}

\begin{thm}
\label[thm]{theorem: characterizations of hypocentral homomorphisms}
Let $f\colon P\twoheadrightarrow Q$ be a surjective quandle homomorphism. Then, the following conditions are equivalent.
\begin{enumerate}[label={(\roman*)}]
    \item\label{item: f is hypocentral}
    The homomorphism $f$ is hypocentral, that is, $\relInn(f)$ is $\Inn(P)$-hypocentral.
    \item\label{item: relative transvection is hypocentral}
    The group $\relTrans(f)$ is $\Inn(P)$-hypocentral.
    \item\label{item: relative adjoint group is hypocentral}
    The group $\relAs(f)$ is $\As(P)$-hypocentral.
    \item\label{item: no perfect subgroup in relative inner group}
    The group $\relInn(f)$ contains no non-trivial $\Inn(P)$-perfect normal subgroup of $\Inn(P)$.
    %\item\label{item: no strongly connected quotient factor} For every quotient factorization $f=h\circ e$ with $e$ strongly connected, the homomorphism $e$ is an isomorphism.
\end{enumerate}
\end{thm}

\begin{proof}
The equivalence of \ref{item: f is hypocentral} and \ref{item: relative transvection is hypocentral} follows from \cref{corollary: relative hypocenter and hypocentral lengths}.
The equivalence of \ref{item: relative transvection is hypocentral} and \ref{item: relative adjoint group is hypocentral} follows from \cref{lemma: hypocentrality under central extensions,proposition: relative adjoint group central extension}.
By \cref{corollary:relative_hypocentral_iff_no_non-trivial_relative_perfect}, \ref{item: f is hypocentral} is equivalent to \ref{item: no perfect subgroup in relative inner group}. 
%It remains to show \ref{item: no perfect subgroup in relative inner group} $\Leftrightarrow$ \ref{item: no strongly connected quotient factor}.
%
%Assume \ref{item: no perfect subgroup in relative inner group} and write $f=h\circ e$, where $e$ is strongly connected. Then $\relTrans(e)$ is $\Inn(P)$-perfect by \cref{theorem: characterization of strongly connected via perfect subgroup} and is contained in $\relInn(f)$. Hence, the assumption yields $\relTrans(e)=1$. Hence, as $e$ is strongly connected, it follows from \cref{proposition: strongly connected and nilpotent is iso} that $e$ is an isomorphism.
%Conversely, let $K\trianglelefteq\Inn(P)$ be a non-trivial $\Inn(P)$-perfect subgroup contained in $\relInn(f)$. Then $K=[\Inn(P),K]\subseteq[\Inn(P),\relInn(f)]\subseteq\relTrans(f)$. Thus, $f$ factors through $P\twoheadrightarrow P/K$, and this quotient is strongly connected by \cref{theorem: characterization of strongly connected via perfect subgroup}. Hence, \ref{item: no strongly connected quotient factor} fails.
\end{proof}

The following generalizes \cref{proposition: strongly connected and nilpotent is iso}.

\begin{prop}
\label[prop]{theorem: strongly connected and hypocentral is iso}
    Any strongly connected and hypocentral homomorphism is an isomorphism.
\end{prop}

\begin{proof}
Suppose that $f$ is both strongly connected and hypocentral. Then, the group $\relTrans(f)$ is $\Inn(P)$-perfect by \cref{proposition:relTrans_of_strongly_conncted_is_Inn-perfect} and is contained in $\relInn(f)$. Hence, the hypocentrality of $f$ yields $\relTrans(f)=1$ by \cref{theorem: characterizations of hypocentral homomorphisms}. Thus, as $f$ is strongly connected and hence is isomorphic to $\pi_{\relTrans(f)}$, it follows that $f$ is an isomorphism.
\end{proof}


\begin{prop}
\label[prop]{proposition: hypocentrality under composition}
    Let $f\colon P\twoheadrightarrow Q$ and $g\colon Q\twoheadrightarrow R$ be surjective quandle homomorphisms. Then, the following hold.
    \begin{enumerate}
        \item\label{item: hypocentrality is closed under composition}
        If $f$ and $g$ are hypocentral, then so is $g\circ f$.
        \item\label{item: hypocentrality descends to the first factor}
        If $g\circ f$ is hypocentral, then $f$ is hypocentral.
    \end{enumerate}
\end{prop}

\begin{proof}
\eqref{item: hypocentrality is closed under composition}:
Put $h=g\circ f$ and let $K\trianglelefteq\Inn(P)$ be an $\Inn(P)$-perfect subgroup contained in $\relInn(h)$. Then $f_*(K)$ is an $\Inn(Q)$-perfect subgroup contained in $\relInn(g)$, as $[\Inn(Q),f_*(K)]=f_*([\Inn(P),K])=f_*(K)$ and $f_*(\relInn(h)) \subseteq \relInn(g)$. The hypocentrality of $g$ forces $f_*(K)=1$. Hence, $K\subseteq\relInn(f)$, and therefore the hypocentrality of $f$ shows $K=1$.
Thus, it follows from \ref{item: no perfect subgroup in relative inner group} of \cref{theorem: characterizations of hypocentral homomorphisms} that $h=g\circ f$ is hypocentral.

\eqref{item: hypocentrality descends to the first factor}: Since $\relInn(f)\subseteq\relInn(g\circ f)$, this also follows from \ref{item: no perfect subgroup in relative inner group} of \cref{theorem: characterizations of hypocentral homomorphisms}.
\end{proof}





\subsection{Strongly connected--hypocentral factorization}
\label{subsection: strongly connected-hypocentral factorization}

In this subsection, we show that every surjection admits a canonical strongly connected--hypocentral factorization and that these two classes form an orthogonal factorization system for surjections in $\Quandle$.

\begin{thm}[Strongly connected--hypocentral factorization]
\label[thm]{theorem: universal strongly connected-hypocentral factorization}
    Let $f\colon P\twoheadrightarrow Q$ be a surjective quandle homomorphism and put $H\coloneqq\hypCenter{f}$. Then, $f$ factors as 
    \[\begin{tikzcd}[column sep=small]
        P \arrow[two heads]{rr}{f} \arrow[two heads]{rd}[swap]{\hypMap{f}} & & Q \\
        & P/H \arrow[two heads]{ru}[swap]{\hypFac{f}}
    \end{tikzcd}\]
    where $\hypMap{f}$ is strongly connected and $\hypFac{f}$ is hypocentral.
\end{thm}

\begin{proof}
The subgroup $H$ is $\Inn(P)$-perfect and is contained in $\relInn(f)$. Hence, $f$ factors through $P/H$ as $f=\hypFac{f}\circ\hypMap{f}$, and $\hypMap{f}$ is strongly connected by \cref{theorem: characterization of strongly connected via perfect subgroup}.

It remains to show that $\hypFac{f}$ is hypocentral. Put $G=\Inn(P)$, $N=\relInn(f)$, and $S=\relInn(\hypMap{f})$. Then, we have $\Inn(P/H) \cong \Inn(P)/\relInn(\hypMap{f}) = G/S$ and, from \cref{lemma: relative symmetry groups under a surjective factorization}, we see $\relInn(\hypFac{f}) \cong \relInn(f)/\relInn(\hypMap{f}) = N/S$. We need to show that $N/S$ is $G/S$-hypocentral.

Since $\hypMap{f}$ is the quotient map by $H$, we have $H \subseteq \relInn(\hypMap{f})=S \subseteq \relInn(f)=N$. We claim that $N/H$ is $G/H$-hypocentral; indeed, if $K/H$ is a $G/H$-perfect subgroup contained in $N/H$, then the condition $[G/H,K/H]=[G,K]/H=K/H$ implies $[G,K] = K$, which shows that $K$ is $G$-perfect. As $H \subseteq K \subseteq N$, the maximality of $H$ forces $H=K$; in other words, $K/H$ is trivial. Thus $N/H$ is $G/H$-hypocentral by \cref{corollary:relative_hypocentral_iff_no_non-trivial_relative_perfect}.

For every $\varphi\in \relInn(\hypMap{f})=S$ and $x\in P$, we have $\hypMap{f}(\varphi(x))=\hypMap{f}(x)$, and hence $\varphi(x)=\xi_x(x)$ for some $\xi_x \in H$. Then,
\[ [\varphi,s_x]=s_{\varphi(x)}s_x^{-1} = s_{\xi_x(x)}s_x^{-1} = [\xi_x,s_x] \in [H,\Inn(P)]=H, \]
and therefore we obtain $[S,G] \subseteq H$. This shows that the induced map $G/H \twoheadrightarrow G/S$ is a central extension.
As $N/H$ is $G/H$-hypocentral, \cref{lemma: hypocentrality under central extensions} yields that $N/S$ is $G/S$-hypocentral, which completes the proof.
\end{proof}


\begin{comment}
\begin{prop} %%% direct proof of orthogonality %%%(ただし未チェック)
\label[prop]{proposition: strongly connected homomorphisms are orthogonal to hypocentral homomorphisms}
    In the category of quandles, strongly connected homomorphisms are orthogonal to hypocentral homomorphisms.
\end{prop}

\begin{proof}
Consider a commutative square
\[
\begin{tikzcd}
    A
    \arrow{r}{u}
    \arrow[two heads]{d}[swap]{e}
    &
    C
    \arrow[two heads]{d}{m}
    \\
    B
    \arrow{r}[swap]{v}
    &
    D,
\end{tikzcd}
\]
where $e$ is strongly connected and $m$ is hypocentral.

Put $R\coloneqq\relAs(e)$, $T\coloneqq\relTrans(e)$, and $L\coloneqq[\Adj(A),R]$. By \cref{proposition: relative adjoint group central extension}, the canonical homomorphism $\rho_A$ maps $R$ onto $T$. Since $e$ is strongly connected, $T$ is $\Inn(A)$-perfect by \cite[Theorem~3.11]{preprint_yuki_tomoki_2026_relativization_of_symmetries_on_quandles}. Hence, $\rho_A(L)=[\Inn(A),T]=T$. Since $\Ker(\rho_A)$ is central, we have $R=L(R\cap\Ker(\rho_A))$, and therefore $[\Adj(A),L]=[\Adj(A),R]=L$.

Put $H\coloneqq\Adj(u)(\Adj(A))$ and $S\coloneqq\Adj(u)(L)$. Then, $S=[H,S]$. Moreover, the commutativity of the square gives $\Adj(m)(S)=\Adj(v)(\Adj(e)(L))=1$, and hence $S\subseteq\relAs(m)$.

Since $m$ is hypocentral, $\relAs(m)$ is $\Adj(C)$-hypocentral by \cref{theorem: characterizations of hypocentral homomorphisms}. Let $\Gamma_\alpha$ denote its transfinite $\Adj(C)$-relative lower central series. We have $S\subseteq\Gamma_0$. If $S\subseteq\Gamma_\alpha$, then $S=[H,S]\subseteq[\Adj(C),\Gamma_\alpha] =\Gamma_{\alpha+1}$, while the limit step is immediate. Thus, transfinite induction gives $S\subseteq\Gamma_\alpha$ for every ordinal $\alpha$. Since this series reaches the trivial subgroup, we obtain $S=1$.

Let $a_1,a_2\in A$ satisfy $e(a_1)=e(a_2)$. Since $T$ acts transitively on every fiber of $e$, there is $t\in T$ such that $t(a_1)=a_2$. Choose $\widetilde t\in L$ with $\rho_A(\widetilde t)=t$. By the naturality of the associated-group action,
\[ u(a_2) = u(\widetilde t\cdot a_1) = \Adj(u)(\widetilde t)\cdot u(a_1) = u(a_1). \]
Hence, $u$ is constant on the fibers of $e$, and therefore induces a unique map $w\colon B\to C$ satisfying $w\circ e=u$.

The map $w$ is a quandle homomorphism: if $b_i=e(a_i)$, then $w(s_{b_1}(b_2)) =u(s_{a_1}(a_2)) =s_{u(a_1)}(u(a_2)) =s_{w(b_1)}(w(b_2))$. Finally, $m\circ w\circ e=m\circ u=v\circ e$, and the surjectivity of $e$ gives $m\circ w=v$. The uniqueness follows from the surjectivity of $e$.
\end{proof}

\end{comment}


%%%% こっちはあってもいいかも: ascendはともかくdescendは直交分解系をなすことからは従わないはず
\begin{comment}
\begin{prop}
\label[prop]{proposition: hypocentral homomorphisms under pullback}
\memo{added prop}
    Consider a pullback square of surjective quandle homomorphisms
\[\begin{tikzcd}
    P'
    \arrow[two heads]{r}{u}
    \arrow[two heads]{d}[swap]{f'}
    &
    P
    \arrow[two heads]{d}{f}
    \\
    R
    \arrow[two heads]{r}[swap]{v}
    &
    Q\rlap{.}
\end{tikzcd}
\]
    Then, for every ordinal $\alpha$, the homomorphism $u_*$ restricts to an isomorphism
    \[ u_*\colon \relTransLCS{\alpha}{f'} \xrightarrow{\sim} \relTransLCS{\alpha}{f}. \]
    In particular, $f'$ is hypocentral if and only if $f$ is hypocentral.
\end{prop}

\begin{proof}
By \cite[Proposition~3.27]{preprint_yuki_tomoki_2026_relativization_of_symmetries_on_quandles}, the restriction of $u_*$ induces an isomorphism
\[ u_*\colon \relTrans(f') \xrightarrow{\sim} \relTrans(f). \]
Moreover, $u\colon P'\twoheadrightarrow P$ is surjective, and hence $u_*\colon\Inn(P')\twoheadrightarrow\Inn(P)$ is surjective.

We prove the assertion by transfinite induction on $\alpha$. The case $\alpha=0$ is the preceding isomorphism. Suppose that the assertion holds for an ordinal $\alpha$. Then,
\begin{align*}
    u_*\bigl(\relTransLCS{\alpha+1}{f'}\bigr)
    &= u_*\bigl( [\Inn(P'),\relTransLCS{\alpha}{f'}] \bigr) \\
    &= [\Inn(P),\relTransLCS{\alpha}{f}] = \relTransLCS{\alpha+1}{f}.
\end{align*}
Since $\relTransLCS{\alpha+1}{f'}\subseteq\relTrans(f')$ and the restriction of $u_*$ to $\relTrans(f')$ is injective, this restriction is an isomorphism in degree $\alpha+1$.

Let $\lambda$ be a limit ordinal, and suppose that the assertion holds for every $\alpha<\lambda$. By definition,
\[ \relTransLCS{\lambda}{f'} = \bigcap_{\alpha<\lambda} \relTransLCS{\alpha}{f'}, \qquad
\relTransLCS{\lambda}{f} = \bigcap_{\alpha<\lambda} \relTransLCS{\alpha}{f}. \]
The inclusion
\[ u_*\left(\bigcap_{\alpha<\lambda} \relTransLCS{\alpha}{f'} \right)
\subseteq \bigcap_{\alpha<\lambda} u_*\bigl(\relTransLCS{\alpha}{f'}\bigr) \]
is immediate.

For the reverse inclusion, let $x\in\bigcap_{\alpha<\lambda}\relTransLCS{\alpha}{f}$. Since $u_*\colon\relTrans(f')\xrightarrow{\sim}\relTrans(f)$ is surjective, there is a unique $x'\in\relTrans(f')$ such that $u_*(x')=x$. For every $\alpha<\lambda$, the induction hypothesis gives
\[ x \in \relTransLCS{\alpha}{f} = u_*\bigl(\relTransLCS{\alpha}{f'}\bigr). \]
Hence there is $x'_\alpha\in\relTransLCS{\alpha}{f'}$ such that $u_*(x'_\alpha)=x$. The injectivity of $u_*$ on $\relTrans(f')$ gives $x'_\alpha=x'$. Thus, $x'\in\relTransLCS{\alpha}{f'}$ for every $\alpha<\lambda$, and therefore
\[ x' \in \bigcap_{\alpha<\lambda} \relTransLCS{\alpha}{f'} = \relTransLCS{\lambda}{f'}. \]
It follows that $u_*(\relTransLCS{\lambda}{f'}) =\relTransLCS{\lambda}{f}$. Injectivity again follows from the injectivity of $u_*$ on $\relTrans(f')$.

Thus, $u_*$ restricts to an isomorphism in every ordinal degree. Consequently, $\relTransLCS{\alpha}{f'}=1$ if and only if $\relTransLCS{\alpha}{f}=1$ for every ordinal $\alpha$. The last assertion follows from \cref{theorem: characterizations of hypocentral homomorphisms}.
\end{proof}
\end{comment}


As in \cite[Definition 3.20]{preprint_yuki_tomoki_2026_relativization_of_symmetries_on_quandles}, we refer to a \emph{pre-covering} homomorphism $f\colon P\to Q$ as a quandle homomorphism satisfying $f(x) = f(y)$ implies $s_x=s_y$ for all $x,y\in P$. Surjective pre-covering homomorphisms are precisely covering homomorphisms.
The following orthogonality result is established in the same reference.


\begin{prop}[{\cite[Proposition 3.21]{preprint_yuki_tomoki_2026_relativization_of_symmetries_on_quandles}}]
\label[prop]{lemma: strongly connected homomorphisms and covering condition}
    The class of strongly connected homomorphisms is left orthogonal to the class of pre-covering homomorphisms in the category of quandles.
\end{prop}

Since (pre-)covering maps are not closed under composition, they cannot form the right class of an orthogonal factorization system. We overcome this obstruction by replacing coverings with hypocentral homomorphisms.

We use transfinite cocompositions in the usual sense: at every nonzero limit ordinal, the corresponding object is the limit of the preceding tower.

\begin{comment}
\begin{lem} %%%% well-known fact %%%%%
\label[lem]{lemma: orthogonality under transfinite cocomposition}
    Let $e\colon A\to B$ be a quandle homomorphism, and let $X_\lambda\to X_0$ be the transfinite cocomposite of an inverse ordinal tower $(X_\alpha)_{\alpha\leq\lambda}$. If $e$ is left orthogonal to $X_{\alpha+1}\to X_\alpha$ for every $\alpha<\lambda$, then $e$ is left orthogonal to $X_\lambda\to X_0$.
\end{lem}

\begin{proof}
Consider a commutative square from $e$ to
$X_\lambda\to X_0$.
We construct compatible homomorphisms
$w_\alpha\colon B\to X_\alpha$ by transfinite induction.

In degree $0$, let $w_0$ be the bottom homomorphism of the square.
Suppose that $w_\alpha$ has been constructed. The lifting property
against $X_{\alpha+1}\to X_\alpha$ gives a unique compatible
homomorphism $w_{\alpha+1}$.
If $\beta$ is a nonzero limit ordinal, the compatible family
$(w_\alpha)_{\alpha<\beta}$ induces a unique homomorphism
$w_\beta\colon B\to X_\beta$ by the universal property of the inverse
limit. Thus, $w_\lambda$ is the required diagonal filler.
The same induction proves uniqueness.
\end{proof}
\end{comment}

\begin{thm}
\label[thm]{theorem: transfinite covering characterization of hypocentral homomorphisms}
    Let $f\colon P\twoheadrightarrow Q$ be a surjective quandle homomorphism. Then, the following conditions are equivalent.
\begin{enumerate}[label={(\roman*)}]
    \item\label{item: homomorphism is hypocentral}
    $f$ is hypocentral.
    \item\label{item: transfinite cocomposite of covering condition}
    $f$ is a transfinite cocomposite of pre-covering homomorphisms.
    \item\label{item: hypocentral right orthogonal}
    $f$ is right orthogonal to every strongly connected homomorphism.
\end{enumerate}
\end{thm}

\begin{proof}
\ref{item: homomorphism is hypocentral} $\Rightarrow$ \ref{item: transfinite cocomposite of covering condition}:
Choose an ordinal $\lambda$ such that $\relInnLCS{\lambda}{f}=1$, and put $C_\alpha\coloneqq P/\relInnLCS{\alpha}{f}$ for $\alpha\leq\lambda$.
Then, $C_\lambda=P$; the induced homomorphism $C_0=P/\relInn(f)\twoheadrightarrow Q$ is rigid and hence a covering homomorphism; and every successor homomorphism $C_{\alpha+1}\twoheadrightarrow C_\alpha$ is a covering by \cref{lemma: quotient covering via commutator}.
So, the induced decomposition tower of $f$
\[\begin{tikzcd}[column sep=small]
P=C_\lambda \arrow[two heads]{r} &
\cdots \arrow[two heads]{r} &
C_{\alpha+1} \arrow[two heads]{r} &
C_{\alpha} \arrow[two heads]{r} &
\cdots \arrow[two heads]{r} &
C_0 \arrow[two heads]{r} & Q
\end{tikzcd}\]
consists of covering homomorphisms at every successor stage. However, at a limit ordinal $\beta$, $C_\beta$ is not necessarily the limit of the preceding tower; so, this tower cannot be regarded as a transfinite cocomposition of covering homomorphisms.

Instead, we refine this tower at its limit stages. Suppose that $\beta\leq\lambda$ is a nonzero limit ordinal, and let $\textstyle L_\beta = \lim_{\alpha<\beta} C_{\alpha}$ be the inverse limit of the part of the refined tower already constructed below $\beta$. The compatible quotient homomorphisms induce a canonical homomorphism $g_\beta\colon C_\beta\to L_\beta$.

We claim that $g_\beta$ satisfies the pre-covering condition. Let $x,y\in P$ and suppose that the classes of $x$ and $y$ in $C_\beta=P/\relInnLCS{\beta}{f}$ have the same image under $g_\beta$. Then, for every ordinal $\alpha<\beta$, their classes in $C_\alpha$ necessarily coincide, and there is $\varphi_\alpha\in\relInnLCS{\alpha}{f}$ such that $y=\varphi_\alpha(x)$. Therefore, $s_ys_x^{-1} = [\varphi_\alpha,s_x] \in \relInnLCS{\alpha+1}{f}$ for every $\alpha<\beta$. Since $\beta$ is a limit ordinal, it follows that $\textstyle s_ys_x^{-1}\in\relInnLCS{\beta}{f}=\bigcap_{\alpha<\beta}\relInnLCS{\alpha}{f}$. Thus, the classes of $x$ and $y$ induce the same structure map on $C_\beta=P/\relInnLCS{\beta}{f}$, proving that $g_\beta$ is pre-covering.

At every nonzero limit ordinal $\beta$, insert $L_\beta$ immediately below $C_\beta$ with successor homomorphism $g_\beta$. Reindexing the resulting tower gives a transfinite cocomposition of homomorphisms satisfying the pre-covering condition, whose cocomposite is $f$.

\ref{item: transfinite cocomposite of covering condition} $\Rightarrow$ \ref{item: hypocentral right orthogonal}:
Each homomorphism occurring in the tower is right orthogonal to every strongly connected homomorphism by \cref{lemma: strongly connected homomorphisms and covering condition}. Since right orthogonality is preserved under transfinite cocompositions, the transfinite cocomposite is also right orthogonal to strongly connected homomorphisms.

\ref{item: hypocentral right orthogonal} $\Rightarrow$ \ref{item: homomorphism is hypocentral}:
Suppose that $f$ is right orthogonal to every strongly connected homomorphism. Considering the commutative diagram
\[\begin{tikzcd}
P \arrow[r, "\id"] \arrow[d, two heads, "\hypMap{f}"]
    & P \arrow[d, two heads, "f"]\\
P/\Hyp(f) \arrow[r, two heads, "\hypFac{f}"']
    & Q\rlap{,}
\end{tikzcd}\]
we have a filler $w\colon P/\Hyp(f) \to P$, since $\hypMap{f}$ is strongly connected (\cref{theorem: universal strongly connected-hypocentral factorization}). Then, $w$ is a retraction of $\hypMap{f}$. Therefore $\hypMap{f}$ is injective, and hence an isomorphism. Thus $f\cong \hypFac{f}$ is hypocentral.
\end{proof}

In particular, every transfinite cocomposite of covering homomorphisms is hypocentral.

\begin{rem}
    A hypocentral homomorphism is a transfinite cocomposite of homomorphisms satisfying the pre-covering condition. However, the homomorphisms inserted at limit stages need not be surjective (see \cref{example: hypocentral but not nilpotent Alexander quandle}). Thus, a hypocentral homomorphism is not necessarily a transfinite cocomposite of covering homomorphisms.
\end{rem}

Even and Gran~\cite[Proposition 3.2]{even_gran_2014_on_factorization_systems_for_surjective_quandle_homomorphisms} previously showed that the classes of connected and rigid homomorphisms form an orthogonal factorization system relative to surjective homomorphisms (see also \cite[Proposition 2.24]{preprint_yuki_tomoki_2026_relativization_of_symmetries_on_quandles}). With the preceding preparations, we conclude that the analogous statement holds for strongly connected homomorphisms.
 
\begin{thm}[Strongly connected--hypocentral factorization system]
\label[thm]{theorem: strongly connected-hypocentral factorization system}
    The pair $(\{\mathrm{strongly\>connected}\}, \{\mathrm{hypocentral}\})$ of classes of surjections is an orthogonal factorization system for the surjections in $\Quandle$.
\end{thm}

\begin{proof}
Every surjective homomorphism admits the factorization in \cref{theorem: universal strongly connected-hypocentral factorization}, and the two classes are orthogonal by \cref{theorem: transfinite covering characterization of hypocentral homomorphisms}.
\end{proof}

\begin{rem}
    From \cite[Proposition 3.32]{preprint_yuki_tomoki_2026_relativization_of_symmetries_on_quandles}, strongly connected homomorphisms are closed under pullbacks along surjective homomorphisms.
    Therefore, this factorization system for the surjections is in fact \textit{stable} (see \cite[\href{https://ncatlab.org/nlab/show/stable+factorization+system}{stable factorization system}]{preprint_nlab_2008_nlab}).
\end{rem}

\begin{cor}
\label[cor]{corollary: hypocentral homomorphisms under pullback}
    Hypocentral homomorphisms are preserved by pullback along arbitrary quandle homomorphisms.
\end{cor}

\begin{proof}
By \cref{theorem: transfinite covering characterization of hypocentral homomorphisms}, hypocentral homomorphisms are right orthogonal to strongly connected homomorphisms in the category of quandles. Right orthogonality is preserved by pullback. Since a pullback of a surjective homomorphism is surjective, the assertion follows from the same theorem.
\end{proof}




\subsection{Finite-stage factorizations and examples}
\label{subsection: finite-stage strongly connected-nilpotent factorizations}

The strongly connected--hypocentral factorization has a nilpotent second homomorphism precisely when the relative transvection series stabilizes at a finite stage.

\begin{prop}
\label[prop]{proposition: criterion for strongly connected-nilpotent factorization}
    Let $f\colon P\twoheadrightarrow Q$ be a surjective quandle homomorphism. Then, the following conditions are equivalent.
\begin{enumerate}[label={(\roman*)}]
    \item\label{item: strongly connected-nilpotent factorization}
        The homomorphism $f$ admits a factorization $f=m\circ e$ where $e$ is strongly connected and $m$ is nilpotent.
    \item\label{item: canonical hypocentral factor is nilpotent}
        The canonical hypocentral factor $\hypFac{f}$ is nilpotent.
    \item\label{item: finite stabilization of relative transvection series}
        We have $\relTransLCS{c}{f}=\relTransLCS{c+1}{f}$ for some finite $c\geq0$.
\end{enumerate}
If these conditions hold, then $\hypCenter{f}=\relTransLCS{c}{f}$ for every sufficiently large integer $c$, and the canonical factorization is $P\twoheadrightarrow P/\relTransLCS{c}{f}\twoheadrightarrow Q$.
\end{prop}

\begin{proof}
\ref{item: finite stabilization of relative transvection series} $\Rightarrow$ \ref{item: canonical hypocentral factor is nilpotent}:
If $\relTransLCS{c}{f}=\relTransLCS{c+1}{f}$ for some integer $c\geq0$, then $\hypCenter{f}=\relTransLCS{c}{f}$. 
Therefore, $\hypFac{f}$ is in fact $\covTrunc{c}{f}\colon P/\relTransLCS{c}{f} \twoheadrightarrow Q$ in \cref{theorem: universal property of covering length truncation}, which is nilpotent by \cref{theorem: main structure theorem for nilpotent homomorphisms}.

\ref{item: canonical hypocentral factor is nilpotent} $\Rightarrow$ \ref{item: strongly connected-nilpotent factorization}: It follows from \cref{theorem: universal strongly connected-hypocentral factorization}.

\ref{item: strongly connected-nilpotent factorization} $\Rightarrow$ \ref{item: finite stabilization of relative transvection series}:
Let $f=m\circ e$ be the factorization. 
Since $e$ is strongly connected, $\relTrans(e)$ is $\Inn(P)$-perfect by \cref{proposition:relTrans_of_strongly_conncted_is_Inn-perfect}.
It is immediate to see $\relTrans(e) \subseteq \relTrans(f)\subseteq \relInn(f)$. Hence, $\relTrans(e) \subseteq \hypCenter{f}$ by \cref{lemma: relative hypocenter}; in particular, $\relTrans(e) \subseteq \relTransLCS{i}{f}$ for every $i\geq0$.

Since $m$ is nilpotent, there is an integer $c\geq0$ such that $\relTransLCS{c}{m}=1$. By \cref{lemma: relative symmetry groups under a surjective factorization} and \cref{proposition: functoriality of relative lower central series}, we have $e_*(\relTransLCS{c}{f})=\relTransLCS{c}{m}=1$, and hence $\relTransLCS{c}{f}\subseteq \relInn(e)$. Therefore, using \cref{lem:commutator_inn_rel_trans_rel}, we have
\[ \relTransLCS{c+1}{f} = [\Inn(P),\relTransLCS{c}{f}] \subseteq [\Inn(P),\relInn(e)] = \relTrans(e). \]
Together with $\relTrans(e)\subseteq \relTransLCS{c+1}{f}$, we obtain $\relTransLCS{c+1}{f}=\relTrans(e)$, which is $\Inn(P)$-perfect. Consequently, it holds that $\relTransLCS{c+1}{f}=[\Inn(P),\relTransLCS{c+1}{f}]=\relTransLCS{c+2}{f}$.
\end{proof}

\begin{cor}
\label[cor]{corollary: descending chain condition and finite-stage factorization}
    Assume that the $\Inn(P)$-normal subgroups of $\relTrans(f)$ satisfy the descending chain condition. Then, $f$ admits a factorization into a strongly connected homomorphism followed by a nilpotent homomorphism. Moreover, $f$ is hypocentral if and only if it is nilpotent. In particular, these assertions hold when $\relTrans(f)$ is finite.
\end{cor}

\begin{proof}
The assumption shows that the relative transvection series stabilizes at a finite stage, so the first assertion follows from \cref{proposition: criterion for strongly connected-nilpotent factorization}. If $f$ is hypocentral, its stable term is trivial, and hence the series vanishes at a finite stage. Thus, $f$ is nilpotent. The converse is trivial.
\end{proof}

\begin{cor}
\label[cor]{corollary: strongly connected-nilpotent factorization system for finite inner groups}
    On the full subcategory of $\Quandle$ consisting of quandles with finite inner automorphism groups, the classes of strongly connected and nilpotent homomorphisms form a (stable) orthogonal factorization system for surjections. In particular, the same holds for finite quandles.
\end{cor}

\begin{proof}
For every homomorphism in this category, the relative transvection group is finite, so hypocentrality and nilpotency coincide by \cref{corollary: descending chain condition and finite-stage factorization}. The canonical factorization remains in the subcategory, since $\Inn(P/\hypCenter{f})$ is a quotient of $\Inn(P)$. Thus, the orthogonal factorization system of \cref{theorem: strongly connected-hypocentral factorization system} restricts to this subcategory, and hypocentrality can be replaced by nilpotency.
(Moreover, the subcategory is closed under pullbacks: the inner automorphism group of a pullback is a quotient of a subgroup of the product of the inner automorphism groups of its two factors. Hence the resulting orthogonal factorization system is stable as well.)
\end{proof}

Finally, we observe that in general, the classes of nilpotent and hypocentral homomorphisms do not coincide. In particular, strongly connected and nilpotent homomorphisms do not form a factorization system on all quandles.


\begin{eg}
\label[eg]{example: hypocentral but not nilpotent Alexander quandle}
    Let $P\coloneqq\Alex(\ZZ,-\id)$, and let $f\colon P\twoheadrightarrow\{*\}$ be the unique homomorphism into the terminal quandle. 
    By \cref{example: relative Alexander lower central series}, we have
    \[
        \relTransLCS{i}{f}=\{L_{2^{i+1}n}\mid n\in\ZZ\} \qquad (i<\omega),
    \]
    and hence $\relTransLCS{\omega}{f}=1$. Thus, $f$ is hypocentral but not nilpotent. 
    Moreover, $f$ is not a transfinite cocomposite of covering homomorphisms.

    Since $s_x(0)=2x$ for every $x\in P$, the quandle $P$ is faithful, that is, $s_x= s_y$ implies $x=y$. Thus, every covering homomorphism with domain $P$ is an isomorphism. Moreover, every non-injective surjective homomorphism $q\colon P\twoheadrightarrow X$ has finite target and infinite fibers. Indeed, if $q(a)=q(b)$ with $a\neq b$, then, for every $z\in\ZZ$, we have
    \[
        q(z+2(a-b))=q(s_as_b^{-1}(z))
        =s_{q(a)}s_{q(b)}^{-1}(q(z))=q(z).
    \]
    Hence, $q$ has a nonzero period, so $|X|\leq2|a-b|$ and every fiber of $q$ is infinite.

    Suppose that $f$ is the cocomposite of a continuous tower $(X_\alpha)_{\alpha\leq\lambda}$ with $X_0=\{*\}$, $X_\lambda=P$, and every successor homomorphism $X_{\alpha+1}\twoheadrightarrow X_\alpha$ covering. Write $q_\alpha\colon P\to X_\alpha$ for the canonical homomorphisms. Successive lifting, using continuity at limit ordinals, shows that every $q_\alpha$ is surjective.

    Let $\beta$ be the least ordinal for which $q_\beta$ is an isomorphism. Since $P\neq\{*\}$, we have $\beta>0$. If $\beta=\alpha+1$, then $q_\alpha$ is a covering with domain $P$, and hence an isomorphism, contradicting the minimality of $\beta$. Thus, $\beta$ is a limit ordinal, and
    \[
        P\cong X_\beta\cong\varprojlim_{\alpha<\beta}X_\alpha.
    \]
    For every $\alpha<\beta$, the quandle $X_\alpha$ is finite and every fiber of $q_\alpha$ is infinite.

    Give each $X_\alpha$ the discrete topology and $P$ the resulting inverse limit topology. Then, $P$ is a nonempty countable compact Hausdorff space. Since the index set is linearly ordered, the fibers of the $q_\alpha$ form a basis of open sets. These fibers are all infinite, so $P$ has no isolated points. This contradicts the Baire category theorem, since every nonempty countable compact Hausdorff space has an isolated point.
\end{eg}
